% from ega1-8-fr.tex
\oldpage[I]{167}
Cela résulte de (\hyperref[I.8.2.3-fr]{8.2.3}) et de l'axiome
$(T_0)$ (\hyperref[I.2.1.4-fr]{2.1.4}).

\begin{corollary}[8.2.5]
\label{I.8.2.5-fr}
Soit $X$ un schéma intègre dont l'espace sous-jacent est noethérien~;
lorsque $f$ parcourt le corps $R$ des fonctions rationnelles sur $X$, les
ensembles $\delta(f)$ engendrent la topologie de $X$.
\end{corollary}

En effet, toute partie fermée de $X$ est alors réunion finie d'ensembles
fermés irréductibles, c'est-à-dire de la forme $\overline{\{y\}}$
(\hyperref[I.2.1.5-fr]{2.1.5}). Or, si
$x\notin\overline{\{y\}}$, il existe une fonction rationnelle $f$ définie en
$x$ et non en $y$ (\hyperref[I.8.2.3-fr]{8.2.3}), autrement dit, on a
$x\in\delta(f)$ et $\delta(f)$ ne rencontre pas $\overline{\{y\}}$. Le
complémentaire de $\overline{\{y\}}$ est par suite réunion d'ensembles de la
forme $\delta(f)$, et en vertu de la première remarque, tout ouvert de $X$ est
réunion d'intersections finies d'ouverts de la forme $\delta(f)$.

\begin{env}[8.2.6]
\label{I.8.2.6-fr}
Le cor. (\hyperref[I.8.2.5-fr]{8.2.5}) montre que la topologie de $X$ est
entièrement caractérisée par la donnée de la famille d'anneaux locaux
$(\mathscr{O}_x)_{x\in X}$ ayant $R$ pour corps des fractions. Il revient au
même d'ailleurs de dire que les parties fermées de $X$ sont définies de la
façon suivante~: étant donnée une partie finie
$\{x_1,\ldots,x_n\}$ de $X$, on considère l'ensemble des $y\in X$ tels que
$\mathscr{O}_y\subset\mathscr{O}_{x_i}$ pour un indice $i$ au moins, et ces
ensembles (pour tous les choix de $\{x_1,\ldots,x_n\}$) sont les ensembles
fermés de $X$. En outre, une fois connue la topologie de $X$, le faisceau
structural $\mathscr{O}_X$ est aussi bien déterminé par la famille des
$\mathscr{O}_x$, puisque
$\Gamma(U,\mathscr{O}_X)=\bigcap_{x\in U}\mathscr{O}_x$
(\hyperref[I.8.2.1.1-fr]{8.2.1.1}). La famille
$(\mathscr{O}_x)_{x\in X}$ détermine donc complètement le préschéma $X$
lorsque $X$ est un schéma intègre dont l'espace sous-jacent est noethérien.
\end{env}

\begin{proposition}[8.2.7]
\label{I.8.2.7-fr}
Soient $X$, $Y$ deux schémas intègres, $f:X\to Y$ un morphisme dominant
(\hyperref[I.2.2.6-fr]{2.2.6}), $K$ (resp. $L$) le corps des fonctions
rationnelles sur $X$ (resp. $Y$). Alors $L$ s'identifie à un sous-corps de
$K$, et pour tout $x\in X$, $\mathscr{O}_{f(x)}$ est l'unique anneau local de
$Y$ dominé par $\mathscr{O}_x$.
\end{proposition}

En effet, si $f=(\psi,\theta)$ et si $a$ est le point générique de $X$,
$\psi(a)$ est le point générique de $Y$
(\hyperref[0.2.1.5-fr]{0, 2.1.5})~; $\theta_a^\sharp$ est par suite un
monomorphisme du corps $L=\mathscr{O}_{\psi(a)}$ dans le corps
$K=\mathscr{O}_a$. Comme tout ouvert affine non vide $U$ de $Y$ contient
$\psi(a)$, il résulte de (\hyperref[I.2.2.4-fr]{2.2.4}) que l'homomorphisme
$\Gamma(U,\mathscr{O}_Y)\to\Gamma(\psi^{-1}(U),\mathscr{O}_X)$ correspondant
à $f$ est la restriction de $\theta_a^\sharp$ à
$\Gamma(U,\mathscr{O}_Y)$. Donc, pour tout $x\in X$,
$\theta_x^\sharp$ est la restriction à $\mathscr{O}_{\psi(x)}$ de
$\theta_a^\sharp$, et est par suite un monomorphisme. On sait en outre que
$\theta_x^\sharp$ est un homomorphisme local, donc, si on identifie $L$ à un
sous-corps de $K$ par $\theta_a^\sharp$, $\mathscr{O}_{\psi(x)}$ est dominé
par $\mathscr{O}_x$ (\hyperref[I.8.1.1-fr]{8.1.1})~; c'est d'ailleurs le seul
anneau local de $Y$ dominé par $\mathscr{O}_x$ puisque deux anneaux locaux de
$Y$ qui sont apparentés sont identiques
(\hyperref[I.8.2.2-fr]{8.2.2}).

\begin{proposition}[8.2.8]
\label{I.8.2.8-fr}
Soient $X$ un préschéma irréductible, $f:X\to Y$ une immersion locale
(resp. un isomorphisme local)~; on suppose en outre le morphisme $f$ séparé.
Alors $f$ est une immersion (resp. une immersion ouverte).
\end{proposition}

Soit $f=(\psi,\theta)$~; il suffit, dans les deux cas, de prouver que $\psi$
est un \emph{homéomorphisme} de $X$ sur $\psi(X)$
(\hyperref[I.4.5.3-fr]{4.5.3}). Remplaçant $f$ par $f_{\mathrm{red}}$
(\hyperref[I.5.1.6-fr]{5.1.6} et
\hyperref[I.5.5.1-fr]{5.5.1, (vi)}), on peut supposer $X$ et $Y$ réduits. Si
$Y'$ est le sous-préschéma fermé réduit de $Y$ ayant pour espace sous-jacent
$\overline{\psi(X)}$, $f$ se factorise en
$X\xrightarrow{f'}Y'\xrightarrow{j}Y$, où $j$ est l'injection canonique
(\hyperref[I.5.2.2-fr]{5.2.2}). Il résulte de
(\hyperref[I.5.5.1-fr]{5.5.1, (v)}) que $f'$ est encore un morphisme séparé~;
en outre, $f'$ est encore

